\section{Argument strength and probabilistic assessment}
\label{sec:p18-r55-proofs}

\begin{annexitem}{annex:p18.r55.detection}{thm:p18-r55-detection}{Positive and negative evidential force}
\begin{proof}
The stipulated joint probabilities are
\[
 \mathbb P(H\cap E)=\pi s,\quad
 \mathbb P(H^c\cap E)=(1-\pi)(1-c),\quad
 \mathbb P(H\cap E^c)=\pi(1-s),\quad
 \mathbb P(H^c\cap E^c)=(1-\pi)c.
\]
They are nonnegative and sum to one for \(s,c\in[0,1]\).
Adding the first two gives
\(D_+=\mathbb P(E)=\pi s+(1-\pi)(1-c)\);
adding the last two gives
\(D_-=\mathbb P(E^c)=\pi(1-s)+(1-\pi)c\).
Dividing the relevant joint probability by its positive evidence
probability is exactly the conditioning definition and gives the
two displayed posteriors. No division is asserted if the respective
denominator is zero. For their changes from prior belief,
\[
 p_+-\pi=\frac{\pi(1-\pi)(s+c-1)}{D_+},\qquad
 q_--(1-\pi)=\frac{\pi(1-\pi)(s+c-1)}{D_-}.
\]
Since \(0<\pi<1\) and the relevant denominator is positive,
the signs of these differences are exactly the sign of \(s+c-1\).
This proves all three comparison cases under the stated conditions.

For the first written numerical example the four masses are
\(8/100,9/100,2/100,81/100\), ordered as
\(H\cap E,H^c\cap E,H\cap E^c,H^c\cap E^c\).
Thus the posteriors are \((8/100)/(17/100)=8/17\)
and \((81/100)/(83/100)=81/83\).
For \(s=1/50,c=49/50\), the probabilities of \(E\)
under \(H\) and \(H^c\) are both \(1/50\).
Likewise their absence probabilities are both \(49/50\).
Both updates therefore leave the respective prior unchanged.
These calculations verify the example's distinction between high
confidence and an increase supplied by the particular observation.
\end{proof}
\end{annexitem}

\begin{annexitem}{annex:p18.r55.norm-boundaries}{thm:p18-r55-norm-boundaries}{Four distinct conditional standards}
\begin{proof}
If \(E\subseteq H\), then \(E\cap H=E\), so positive-probability
conditioning gives \(\mathbb P(H\mid E)=1\). For the failure
of the converse, take \(\Omega=\{x,y\}\), masses one and zero,
\(E=\Omega\) and \(H=\{x\}\). The posterior is one,
but \(y\in E\setminus H\). The event implication is false
at that probability-zero outcome. A logical entailment conclusion
would require ruling out such countercases, not just assigning
them zero probability.

For a binary truth indicator with probability \(p\) of one,
the expected quadratic loss of report \(r\) is
\[
 p(r-1)^2+(1-p)r^2
  =r^2-2pr+p=(r-p)^2+p(1-p).
\]
The squared term is nonnegative and vanishes exactly at \(r=p\),
which lies in the allowed interval. This proves the unique minimizer,
including \(p=0\) and \(p=1\). When \(p\) is obtained by
conditioning on the actual evidence, the calculation concerns that
conditional distribution. It assumes the indicated loss, not a
universal preference ordering over all epistemic and social outcomes.

For repeated identical reports, \(E_1\cap E_2=E_1\).
The two conditional probabilities are consequently the same defined
ratio. No assumption of independence is used; identical nontrivial
events are not independent. Merely supplying the same evidence in
another sentence does not alter this equality.

Finally the proposed action's expected utility is
\(10(1-q)-100q=10-110q\), while the alternative's is zero.
Their strict comparison is therefore \(q<1/11\); equality gives
the stated tie and larger \(q\) favors the alternative. This
calculation depends on the stipulated consequences, utilities and
expected-utility choice convention. A different loss or feasible
alternative need not have this threshold.
\end{proof}
\end{annexitem}

\begin{annexitem}{annex:p18.r55.graph-boundary}{thm:p18-r55-graph-boundary}{Extension acceptance does not identify probabilistic force}
\begin{proof}
In the supplied attack-free framework, every set defends \(a\),
so its characteristic operator sends both \(\varnothing\)
and \(\{a\}\) to \(\{a\}\). Its sole conflict-free
fixed point is \(\{a\}\), proving complete and grounded
uniqueness. Both subsets are admissible, and only \(\{a\}\)
is maximal, proving preferred uniqueness. The empty set cannot
attack the outside argument; \(\{a\}\) has no outside argument
to attack. This proves stable uniqueness as well, directly from
the existing definitions.

For the probability comparison take four outcomes representing
\(H\cap E,H\cap E^c,H^c\cap E,H^c\cap E^c\).
The first probability model gives them masses
\(3/8,1/8,1/8,3/8\); the second gives them
\(1/8,3/8,3/8,1/8\). Every mass is positive and each list
sums to one. In both models \(\mathbb P(H)=1/2\) and
\(\mathbb P(E)=1/2\). Their conditional probabilities of
\(H\) given \(E\) are respectively
\((3/8)/(1/2)=3/4\) and \((1/8)/(1/2)=1/4\).
They lie above and below their common prior, respectively.

The attack graph and its four extension families are unchanged.
They contain none of those likelihood or mass assignments as input.
Thus a function of this graph and its extension verdict alone cannot
identify which of the two opposite evidential comparisons holds.
An additional rule linking probabilistic evidence to graph construction
could impose extra restrictions, but it was not among the supplied
abstract inputs and would need a separate specification and justification.
\end{proof}
\end{annexitem}
